\documentclass[11pt]{article}
\usepackage[margin=1in]{geometry}
\usepackage{amsmath,amssymb,amsthm,mathtools}
\usepackage[T1]{fontenc}
\usepackage{lmodern,microtype,enumitem,xurl}
\usepackage[colorlinks=true,linkcolor=blue,citecolor=blue,urlcolor=blue]{hyperref}
\newtheorem{theorem}{Theorem}[section]
\newtheorem{proposition}[theorem]{Proposition}
\newtheorem{corollary}[theorem]{Corollary}
\newtheorem{lemma}[theorem]{Lemma}
\theoremstyle{remark}\newtheorem{remark}[theorem]{Remark}
\newcommand{\R}{\mathbb R}
\newcommand{\E}{\mathbb E}
\newcommand{\F}{\mathrm F}
\newcommand{\IP}{\operatorname{IP}}
\newcommand{\TC}{\mathsf{TC}}
\newcommand{\norm}[1]{\lVert #1\rVert}
\title{The cost of narrow input layers in Boolean neural networks}
\author{ChatGPT-6 Astra \and Kaarel H\"anni \and Dmitry Vaintrob}
\date{Research draft --- 25 October 2026}
\hypersetup{pdftitle={The cost of narrow input layers in Boolean neural networks},pdfauthor={ChatGPT-6 Astra; Kaarel Haenni; Dmitry Vaintrob}}
\begin{document}
\maketitle
\begin{abstract}
A narrow neural network can represent a simple Boolean function by storing
its inputs at extremely fine resolution. We quantify the cost of this
bottleneck. If the first affine layer has rank $r$ and the product of layer
Frobenius norms is $P$, a globally unit-margined Boolean function has a
constant-depth threshold circuit of size $(d+\max(1,P)+r+2)^{O(r+1)}$.
For inner-product parity on two $n$-bit blocks, every representation whose
first layer has width $r$ requires $P\ge(2^{(n-1)/r}-1)/\sqrt2$, regardless
of its subsequent depth and widths. A first layer on a common rational
grid must also have exponentially fine resolution when $r$ is fixed.
Yet the same function has a depth-three tanh representation of width
$O(n)$ and product $O(n^{5/2}\log^2(n+2))$. These results isolate a
representation obstruction: expanding the first layer can eliminate an
exponential norm cost. They do not establish a learning lower bound for
an unrestricted Bayesian neural-network prior.
\end{abstract}

\section{Model and purpose}
Let $x\in\{-1,1\}^d$ and $\bar x=(1,x)$. Consider a classifier
\begin{equation}\label{eq:model}
 z(x)=G(W\bar x),\qquad f(x)z(x)\ge1\quad\text{for every }x,
\end{equation}
where $f$ takes values in $\{-1,1\}$, $W$ is a real rectangular matrix,
and $G$ is $T$-Lipschitz in Euclidean norm. The column of $W$ multiplying the augmented coordinate $1$
is the actual charged bias column. Write
\[
 F=\norm W_\F,\qquad A=TF,\qquad r=\operatorname{rank}W.
\]
For an ordinary neural network with neuronwise $1$-Lipschitz activations,
$T$ may be taken to be the product of the remaining layer Frobenius norms,
including the linear readout. The first activation is included in $G$.
Thus $A$ can be taken equal to the full layer product $P$.
No assumption about efficient evaluation, zero preservation, or
stability of the activations is needed for the first two theorems below.

A small, deep teacher may have a cheaper alternative implementation.
The results here explain why preserving its narrow first layer can be
expensive. They concern representations, not posterior probabilities.

A threshold gate tests an affine inequality in Boolean inputs. The class
$\TC^0$ consists of nonuniform families of polynomial-size, constant-depth
circuits with majority and negation gates. Polynomially bounded integer
weights can be replaced by repeated inputs and constant wires, with
polynomial size overhead. All complexity statements below are nonuniform:
the circuit may hardcode a finite table of labels.

The proof first rounds a low-rank code using its global margin. For
inner-product parity, distinct truth-table rows then force separated codes,
giving norm and common-grid precision bounds. An explicit wider network
shows how expanding the first layer removes these costs.

\section{A low-rank first layer gives a constant-depth circuit}
\begin{theorem}\label{thm:rank}
Every classifier satisfying \eqref{eq:model} has a constant-depth
threshold circuit of size
\begin{equation}\label{eq:size}
 \bigl(d+\max(1,A)+r+2\bigr)^{O(r+1)}.
\end{equation}
The implicit constant is universal. Its threshold gates can be expanded
into unweighted majority gates within the same size bound. In particular,
fixed first-layer rank and polynomial Frobenius product imply $\TC^0$,
independently of all later widths and of network depth.
\end{theorem}
\begin{proof}
If $r=0$, or if $T=0$, the function is constant. Otherwise take a singular
value factorization $W=UB$, where $U:\R^r\to\operatorname{im}W$ is an
isometry and $\norm B_\F=F>0$. Set
\[
 u(x)=B\bar x/F,\qquad g(u)=G(FUu).
\]
Then $\norm{u(x)}\le R:=\sqrt{d+1}$, each affine coefficient of $u$ has
absolute value at most one, and $g$ is $A$-Lipschitz. Oppositely labelled
inputs therefore have code separation at least $2/A$.

Put $M=\max(1,A)$. Choose dyadic $h$ such that
\[
 \frac1{32M\sqrt r}\le h\le\frac1{16M\sqrt r}.
\]
Choose another dyadic spacing $\rho\le h/[8(d+1)]$, with $h/\rho$ a power
of two and $1/\rho$ polynomially bounded in $d+M+r$. Round every affine
coefficient of $u$ to this grid. The resulting code $\widetilde u(x)$
differs in each coordinate from $u(x)$ by at most $h/8$.

Round each coordinate of $\widetilde u(x)$ to a nearest multiple of $h$,
with a fixed tie rule. If two inputs receive the same integer grid code,
their exact codes differ in each coordinate by at most $h+2(h/8)$.
Their Euclidean distance is consequently at most
\[
 \frac54\sqrt r\,h<\frac2A.
\]
They have the same label. Each integer code coordinate has magnitude at
most $C(1+M\sqrt{r(d+1)})$, so the number of possible codes is bounded by
$(d+M+r+2)^{O(r)}$.

Hardcode the label of every attained code, assigning arbitrary labels to
unattained ones. Membership in a grid cell is a conjunction of $2r$
affine threshold tests on the original input. An OR over the cells labelled
positive computes $f$ using three threshold layers. Half-open endpoints
or exact ties are implemented by the chosen integer comparison convention.
Multiplication by the common denominator makes every test an integer
threshold with weights and thresholds polynomially bounded in $d+M+r$.
Expanding those weights into repeated literals, and implementing AND and
OR by majority gates with constants, gives \eqref{eq:size}.
\end{proof}

\begin{remark}
The theorem does not require that the original first-layer coefficients
have bounded precision. Its safe rounding scale follows from the product
bound and the global margin. If $P$ is exponential in depth, its conclusion
can still be exponential in depth. Thus it does not by itself establish a
polynomial-depth learning bound for arbitrary narrow teachers.
\end{remark}

\section{Inner-product parity requires information through the first layer}
For $n\ge1$ and $a,b\in\{-1,1\}^n$, define
\[
 \IP_n(a,b)=(-1)^{\sum_{i=1}^n((1-a_i)/2)((1-b_i)/2)}.
\]
This function has a Boolean circuit with $O(n)$ bounded-fan-in gates:
compute the $n$ pairwise conjunctions, then their parity.

\begin{theorem}[Product lower bound]\label{thm:packing}
Suppose
\[
 z(a,b)=G(Aa+Bb+c),\qquad \IP_n(a,b)z(a,b)\ge1
\]
on the full Boolean cube, and $G$ is $T$-Lipschitz. Let
$r_A=\operatorname{rank}A$. Then $r_A\ge1$ and
\begin{equation}\label{eq:packing}
 T\norm A_\F\ge\frac{2^{(n-1)/r_A}-1}{\sqrt2}.
\end{equation}
In a $1$-Lipschitz-activation network with first hidden width at most $r$,
the full Frobenius product therefore obeys
\begin{equation}\label{eq:product}
 P\ge\frac{2^{(n-1)/r}-1}{\sqrt2}.
\end{equation}
Subsequent widths and depth are unrestricted.
\end{theorem}
\begin{proof}
For distinct $a,a'$, select a coordinate $i$ where they differ and choose
$b$ with $b_i=-1$ and every other coordinate $+1$. Then
$\IP_n(a,b)=a_i$ and $\IP_n(a',b)=a'_i$. The outputs differ by at least
$2$, while Lipschitz continuity bounds their difference by
$T\norm{Aa-Aa'}$. Hence
\begin{equation}\label{eq:separation}
 \norm{Aa-Aa'}\ge2/T
\end{equation}
for all distinct $a,a'$. In particular, $T>0$ and $r_A\ge1$.

For uniform $a$ on the cube, independent symmetric signs give
\[
 \E\norm{Aa}^2=\norm A_\F^2.
\]
At least $2^{n-1}$ code points therefore lie in the $r_A$-dimensional
image ball of radius $\sqrt2\norm A_\F$. Place a radius-$1/T$ ball around
each. Their interiors are disjoint by \eqref{eq:separation}, and they
are contained in the ball of radius $\sqrt2\norm A_\F+1/T$. Comparing
$r_A$-dimensional volumes gives
\[
 2^{n-1}\le(1+\sqrt2 T\norm A_\F)^{r_A},
\]
which is \eqref{eq:packing}. In the neural-network specialization,
$T\norm A_\F\le P$ and $r_A\le r$, proving \eqref{eq:product}.
\end{proof}

\begin{corollary}
For fixed first hidden width, every unit-margin representation of
$\IP_n$ has product $\exp(\Omega(n))$. If $P$ is polynomial in $n$,
the first layer must instead have rank $\Omega(n/\log(n+2))$ on the
$a$ input block.
\end{corollary}

The next obstruction concerns input resolution rather than norm. It
requires no continuity or Lipschitz hypothesis at all.
\begin{theorem}[Common-grid precision lower bound]\label{thm:grid}
Suppose a deterministic network computes $\IP_n$ exactly and has first
hidden width $r\ge1$. Let $Q\ge1$ be an integer and $B_0>0$. If every
first-layer coefficient is an integer multiple of $1/Q$ and has magnitude
at most $B_0$, then
\begin{equation}\label{eq:grid}
 2^n\le(2\lfloor QB_0n\rfloor+1)^r,
 \qquad
 Q\ge\frac{2^{n/r}-1}{2B_0n}.
\end{equation}
The later computation may have arbitrary depth, width, activation,
and precision.
\end{theorem}
\begin{proof}
Write the first affine map as $Aa+Bb+c$. The same distinguishing-row
argument shows that $Aa\ne Aa'$ for distinct $a,a'$: otherwise all
subsequent computations would agree for every $b$. Thus $Aa$ has $2^n$
distinct values. Each coordinate of $QAa$ is an integer in
$[-QB_0n,QB_0n]$. Counting those tuples proves \eqref{eq:grid}.
\end{proof}

\begin{remark}
The argument is the elementary deterministic one-way communication
lower bound for inner-product parity. Alice's affine code must distinguish
every possible row of the truth table. Plain parity would not give this
obstruction, because one integer-weight sum retains its parity.

A common denominator is essential to the formulation. Bounding individual
rational denominators separately need not bound their least common multiple
polynomially. For ordinary dyadic coefficients, a polynomial common grid
means $O(\log n)$ fractional bits, whereas \eqref{eq:grid} requires
$\Omega(n/r)$ fractional bits up to logarithmic terms.
\end{remark}

\section{A wider first layer removes the exponential cost}
\begin{theorem}\label{thm:wide}
Inner-product parity has a tanh representation with two hidden layers,
common hidden width $O(n)$, global signed margin greater than one, and
Frobenius product
\[
 P=O\bigl(n^{5/2}\log^2(n+2)\bigr).
\]
The network uses a single augmented constant input and no separately
supplied internal biases. All constant rails and readout coefficients
are charged in the stated product.
\end{theorem}
\begin{proof}
Write $q_i=(1-a_i)(1-b_i)/4\in\{0,1\}$ and
$m=\sum_iq_i$. Let $K=C\log(n+2)$ for a sufficiently large universal
constant $C$. The first hidden layer contains
\[
 h_i=\tanh\bigl(K(-a_i-b_i-1)\bigr),\quad 1\le i\le n,
 \qquad c=\tanh(1).
\]
The ideal value of $h_i$ is $2q_i-1$. Since its signed preactivation
has magnitude at least $K$, the elementary estimate
$1-\tanh u\le2e^{-2u}$ for $u\ge0$ makes its error at most
$1/(100n^2)$. The squared Frobenius norm of this layer is $O(nK^2)$.

Set $s=\sum_i h_i$. For $1\le j\le n$, use second-layer neurons
\[
 t_j=\tanh\bigl(K(s-(-n+2j-1))\bigr).
\]
The constant term is supplied by coefficient
$-K(-n+2j-1)/c$ on the charged rail. Carry the rail itself through
this layer by coefficient $1/c$, so that its new value remains exactly
$c$. Because $s$ is within $1/(100n)$ of $-n+2m$, each threshold has
signed preactivation magnitude at least $K(1-1/(100n))$. Increasing $C$
if necessary makes $t_j$ lie within $1/(100n^2)$ of $+1$ for $j\le m$
and of $-1$ for $j>m$.

Use the linear readout
\[
 z=2\left[1+\sum_{j=1}^n(-1)^j(1+t_j)\right],
\]
with its constant term supplied by the carried $c$ rail. At the ideal
sign values this equals $2(-1)^m=2\IP_n(a,b)$. Its error is at most
$2n/(100n^2)<1/10$, so its signed margin is greater than one.

The second matrix has squared Frobenius norm $O(n^3K^2)$, including
the constant-shift coefficients. The readout has norm $O(\sqrt n)$:
its coefficient on the constant rail is bounded uniformly, since
$1+\sum_{j=1}^n(-1)^j$ is either zero or one. Multiplying the three
norm bounds gives $O(n^{5/2}K^2)$. Zero padding makes the two hidden
widths equal without changing the represented function or product.
\end{proof}

\section{Consequences for deep narrow teachers and Bayesian learning}
The finite-horizon tanh compiler in \cite[constant-width compiler
appendix]{parallel} gives another representation of the same target:
width $24$, every physical coefficient at most one, and depth polynomial
in $n$. It compiles any $S$-gate Boolean circuit on $d$ inputs to depth
at most $5000(d+S+1)^3$ with unit output margin. Its memory uses
exponentially fine real registers; polynomially many description bits
suffice. This representation fact is the only imported result in this
section; Theorems~\ref{thm:rank}--\ref{thm:wide} have self-contained
proofs above.

Combining that compiler with the present results gives three compatible
facts about inner-product parity:
\begin{enumerate}[leftmargin=*]
\item It has fixed-width, bounded-coefficient, polynomial-depth tanh
teachers.
\item Every unit-margin representation retaining a fixed first hidden
width has exponential Frobenius product, regardless of additional depth.
At bounded coefficients its first layer also needs an exponential common
denominator.
\item Expanding the first hidden layer to width $O(n)$ permits a
polynomial-product, depth-three representation.
\end{enumerate}
The fine resolution of a narrow teacher is therefore a genuine information
requirement, while its large product need not be intrinsic to the function.
A universal economical compiler cannot preserve fixed first-layer width.
An unrestricted learner is free to expand it.

These representation lower bounds do not prove bad Bayesian
generalization. At fixed learner width, ordinary finite-architecture
consistency bounds can give polynomial sample guarantees despite an
exponential product. At large learner width, the inexpensive wide
implementation above is available. A lower bound on one teacher's product,
or even on every representation inside a narrow architecture, cannot be
substituted for a target-relative lower bound on the full posterior.

The unrestricted question---whether some small Boolean teachers require
exponential-in-depth samples under the original Gaussian Bayesian
learner---remains open. The present results explain why faithfully
compressing the given teacher, controlling its perturbations, or merely
counting its first-layer cells can miss a cheaper implementation of its
Boolean function.

\paragraph{Verification scope.}
The statements in this note have been independently checked at the level
of conventional mathematical proofs. They are not presented as Lean
formalizations. The related analytic-compression formalization in the
companion project proves other geometric and sampling lemmas, not the
complete theorems in this note.

\begin{thebibliography}{9}
\bibitem{parallel}
ChatGPT-6 Astra, Kaarel H\"anni, and Dmitry Vaintrob.
\emph{Frobenius product, parallel circuits, and sample-complexity barriers
for deep Gaussian Bayesian networks}. Research draft, 25 October 2026.
\end{thebibliography}
\end{document}
